\section*{第 \ref{ch:b2:schrodinger-wave-functions} 章——薛定谔方程与波函数}

\begin{solution}{exo:b2:schrodinger-wave-functions:1}
$h/\sqrt{2mE}$：\qty{0.123}{nm}、\qty{12.3}{pm}、\qty{3.9}{pm}（在 \qty{100}{keV} 时，
动能为 $mc^2$ 的五分之一；相对论动量给出
\qty{3.7}{pm}）；中子 \qty{0.18}{nm}；C$_{60}$ \qty{2.8}{pm}；网球
$\sim 10^{-34}\,\unit{m}$。除网球外，所有粒子均已表现出干涉。
\end{solution}

\begin{solution}{exo:b2:schrodinger-wave-functions:2}
$A = (\pi a^2)^{-1/4}$；$\langle x\rangle = 0$，$\Delta x = a/\sqrt2$；$0.84$；$0.9995$。
\end{solution}

\begin{solution}{exo:b2:schrodinger-wave-functions:3}
代入可得 $E\varphi = -(\hbar^2/2m)\varphi'' + V\varphi$。$E/h =
\qty{2.4e14}{Hz}$；\qty{4.8e14}{Hz}、\qty{620}{nm}：即
跃迁所发出的红光。
\end{solution}

\begin{solution}{exo:b2:schrodinger-wave-functions:4}
电子，\qty{1}{nm}：$\tau = \qty{1.7e-14}{s}$，经 \qty{1}{ns} 后为 \qty{60}{\micro m}，
经 \qty{1}{s} 后为 \qty{60}{km}；\qty{1}{mm}：$\tau = \qty{17}{s}$，经 \qty{1}{ns} 后不变，
经 \qty{1}{s} 后变化 $+0.2\%$；质子，\qty{1}{nm}：$\tau = \qty{3.2e-11}{s}$，\qty{30}{nm}，
\qty{30}{m}；尘粒：$\tau = 600$ 年。对于可见物体，$\tau$
天文量级——而其周围环境会不断使它们定域化。
\end{solution}

\begin{solution}{exo:b2:schrodinger-wave-functions:5}
(a) \cref{prop:b2:schrodinger-wave-functions:current}。(b) $|A|^2\hbar k/m$；
$0$；$(|A|^2 - |B|^2)\hbar k/m$。(c) 入射与反射通量
相减，无干涉项：透射系数与反射系数
可由通量定义。(d) $\psi^*\psi'$ 为实数，其
虚部为零。
\end{solution}

\begin{solution}{exo:b2:schrodinger-wave-functions:6}
(a) $A = (2\pi\sigma^2)^{-1/4}$，$\langle x\rangle = 0$，$\Delta x = \sigma$。(b) $\hbar k_0$。
(c) $\sigma \times \hbar/2\sigma = \hbar/2$。(d) $\dd\langle x\rangle/\dd t = \int x\,
\partial_t\rho = -\int x\,\partial_xj = \int j = (\hbar/m)\operatorname{Im}\int\psi^*
\psi' = \langle p\rangle/m$。
\end{solution}

\begin{solution}{exo:b2:schrodinger-wave-functions:7}
(a) \qty{5.9e6}{m/s}、\qty{0.123}{nm}、\qty{51}{ns}。(b) $v_\varphi = \qty{3.0e6}{m/s}$，
$v_{\text{g}} = \qty{5.9e6}{m/s}$：后者。(c) $\tau = \qty{1.7}{\micro s} \gg
\qty{51}{ns}$：不变。(d) $\lambda/a = \qty{4e-7}{rad}$，\qty{0.1}{\micro m}：
几何光学近似成立。
\end{solution}

\begin{solution}{exo:b2:schrodinger-wave-functions:8}
(a) $|\psi|^2$ 及每个期望值均不变。(b) $\tfrac12(|\psi_1|^2
+ |\psi_2|^2) + \operatorname{Re}(\eu^{\iu\alpha}\psi_1^*\psi_2)$。(c) 来自
双缝的波；其相对相位为 $k(r_2 - r_1)$。(d) 探测器
对两条路径处于不同的状态，两个分量
不再干涉，图样是两个单缝的
图样之和。
\end{solution}

\begin{solution}{exo:b2:schrodinger-wave-functions:9}
(a) $\langle V'\rangle = m\omega^2\langle x\rangle$：精确成立。(b) $\langle V'\rangle
= mg$：精确成立。(c) 当波包的宽度远小于
$V''$ 变化的尺度时。(d) 其在 \qty{0.1}{m/s} 下的波长为 $10^{-29}\,\unit{m}$；
只有波包宽度与碗的尺度相当时，此效应才显著。原子的
热波长 \qty{0.1}{nm}，与原子自身尺寸相当。
\end{solution}

\begin{solution}{exo:b2:schrodinger-wave-functions:10}
(a) $\tfrac12(\varphi_1^2 + \varphi_2^2) + \varphi_1\varphi_2\cos\omega_{21}t$。(b) $L/2 +
\cos\omega_{21}t\int x\varphi_1\varphi_2 = L/2 - (16L/9\pi^2)\cos\omega_{21}t$。(c)
$3h^2/8mL^2 = \qty{1.13}{eV}$、\qty{2.7e14}{Hz}、\qty{1.1}{\micro m}。(d) 它以
频率 $\nu_{21}$ 辐射，即谱线；纯定态无
振荡偶极子，因而不辐射。
\end{solution}

\begin{solution}{exo:b2:schrodinger-wave-functions:11}
(a) $\psi \propto \int\eu^{-\sigma_0^2k^2}\eu^{\iu(kx - \hbar k^2t/2m)}\dd k$。(b) $\alpha =
\sigma_0^2 + \iu\hbar t/2m$，$\beta = \iu x$：$|\psi|^2 \propto \exp[-x^2/2\sigma_0^2(1 +
(\hbar t/2m\sigma_0^2)^2)]$。(c) $\Delta v = \hbar/2m\sigma_0$，且 $\Delta v\,\tau =
\sigma_0$。(d) $\omega = ck$：无色散，所有分量速率均为 $c$。
\end{solution}

\begin{solution}{exo:b2:schrodinger-wave-functions:12}
(a) $\Gamma = \hbar/\tau$：\qty{6.6e-8}{eV}；\qty{6.6e-8}{eV}；\qty{66}{MeV}。(b)
$1/2\pi\tau = \qty{16}{MHz}$，比多普勒宽度小一百倍。(c) 在
冷原子与囚禁离子中，多普勒宽度消失，
可达到自然宽度——这是频率标准的基础。(d) $c\tau =
\qty{3}{m}$，即火车的长度。
\end{solution}

\begin{solution}{pb:b2:schrodinger-wave-functions:1}
\textbf{1.} $v = \sqrt{2eU/m} = \qty{8.4e7}{m/s}$；$p = \qty{7.6e-23}{kg.m/s}$；$\lambda =
\qty{8.7}{pm}$。

\textbf{2.} 误差为百分之几：可用于估算（$\lambda =
\qty{8.6}{pm}$ 为精确值）。

\textbf{3.} \qty{4.8}{ns}。

\textbf{4.} $\lambda/a = \qty{3e-8}{rad}$：\qty{12}{nm}——比
一个像素小 $10^4$ 倍。

\textbf{5.} $\Delta\theta = \Delta p_y/p \sim \hbar/2ap = \lambda/4\pi a$：同一量级。

\textbf{6.} $\tau = \qty{1.6}{ms} \gg \qty{4.8}{ns}$：不变。

\textbf{7.} 对于均匀场，$\dd\langle p_y\rangle/\dd t = eE$ 精确成立：
$\langle y\rangle$ 遵循经典抛物线；该管是一种
经典仪器。

\textbf{8.} $v/2 = \qty{4.2e7}{m/s}$；只有群速度是可观测的；
平面波的相位不携带信号（且依赖于能量的任意
零点选取）。

\textbf{9.} $\Delta p = \hbar/2\sigma_0$，$\Delta v = \hbar/2m\sigma_0$。

\textbf{10.} \qty{1.7e-16}{s}、\qty{1.7e-14}{s}、\qty{1.7e-8}{s}；质子 \qty{3.2e-11}{s}；
病毒 \qty{20}{s}。

\textbf{11.} 它处于定态：势场束缚住了
波包；$|\varphi|^2$ 不移动。

\textbf{12.} $v = \qty{5.9e5}{m/s}$；在 $t = \sqrt3\tau = \qty{2.9e-14}{s}$ 时宽度加倍：
\qty{17}{nm}，约十四个波长。

\textbf{13.} 它们的波函数延展于整块晶体：
动量确定，位置不确定。

\textbf{14.} 分辨率 $\sim\lambda/\alpha = \qty{4}{pm}/0.01 = \qty{0.4}{nm}$：
限制孔径的是磁透镜的像差，而非波长。

\textbf{15.} 将 $\lambda$ 与相关尺度（孔径、
$V$ 的尺度）比较，并将扩散时间与观测时间比较。

\textbf{16.} \qty{0.38}{}、\qty{1.50}{}、\qty{3.38}{eV}；$\nu_{21} = \qty{2.7e14}{Hz}$
（\qty{1.1}{\micro m}），$\nu_{31} = \qty{7.3e14}{Hz}$（\qty{410}{nm}）。

\textbf{17.} 相量的模为一；位于 $x = L/2$ 处。

\textbf{18.} 当 $t = 0$ 时，密度堆积在左半部，
在 $T_{21}/2$ 时堆积在右半部。

\textbf{19.} $\langle x\rangle = L/2 - (16L/9\pi^2)\cos\omega_{21}t$：振幅
为 $0.18L = \qty{0.18}{nm}$。

\textbf{20.} 在 $\nu_{21}$ 处对应谱线；定态没有
振荡偶极子，因而是稳定的；叠加态通过辐射
逐步跃迁到 $\varphi_1$。

\textbf{21.} 不确定：$\langle E\rangle = (E_1 + E_2)/2 = \qty{0.94}{eV}$；一次
测量得到 $E_1$ 或 $E_2$，概率各为二分之一。

\textbf{22.} 随机的是：每个电子落在何处；
不随机的是：由这些落点累积形成的图样 $|\psi|^2$。

\textbf{23.} $j = |A|^2\hbar k/m$：若 $|A|^2 = n$ 为单位长度内的电子数，
则每秒的电子数为 $nv$。

\textbf{24.} 方程对时间为一阶，故 $\psi(x,0)$ 即为完整描述，
决定未来；方程取复数形式，使定能量态具有
唯一的时间依赖 $\eu^{-\iu Et/\hbar}$，并得出 $E = p^2/2m$。

\textbf{25.} $\psi$ 为振幅，其平方即为概率；它
服从薛定谔线性方程；定态带有相量
和固定的密度；自由波包以 $p/m$ 运动，并在
$2m\sigma_0^2/\hbar$ 时间内扩散；埃伦费斯特定理在宏观尺度上重新给出牛顿定律。
\end{solution}
